% Editable publishing source. Rebuild without --refresh-tex to preserve edits.
\documentclass[11pt,a4paper]{article}
% Shared editable typography for all five document editions.
\usepackage{fontspec}
\setmainfont{DejaVu Serif}
\setsansfont{DejaVu Sans}
\setmonofont{DejaVu Sans Mono}[Scale=.86]
\usepackage{polyglossia}
\setdefaultlanguage{russian}
\setotherlanguage{english}
\usepackage[a4paper,margin=18mm,top=21mm,bottom=21mm,headheight=15pt]{geometry}
\usepackage{xcolor}
\definecolor{Accent}{HTML}{185A86}
\definecolor{Muted}{HTML}{526573}
\usepackage{graphicx}
\usepackage{longtable,booktabs,array,calc}
\usepackage{fvextra}
\usepackage{pdflscape}
\usepackage{titlesec}
\usepackage{fancyhdr}
\usepackage{enumitem}
\usepackage{etoolbox}
\usepackage{needspace}
\usepackage{tocloft}
\usepackage[unicode,colorlinks=true,linkcolor=Accent,urlcolor=Accent,bookmarksnumbered=false]{hyperref}
\usepackage{bookmark}
\usepackage{microtype}
\AtBeginDocument{\raggedright}
\setcounter{secnumdepth}{0}
\setcounter{tocdepth}{1}
\setlength{\parindent}{0pt}
\setlength{\parskip}{4pt plus 1pt}
\linespread{1.0}
\setlength{\emergencystretch}{3em}
\setlength{\tabcolsep}{4pt}
\renewcommand{\arraystretch}{1.16}
\AtBeginEnvironment{longtable}{\small}
% Measure a compact table or its header and first two data rows. The output
% remains a breakable longtable. Cap the reservation for unusually tall rows.
\newsavebox{\DocumentTableStart}
\newlength{\DocumentTableSpace}
\newcommand{\KeepTableStart}[2]{%
  \setbox\DocumentTableStart=\vbox{\hsize=\linewidth\small #2}%
  \setlength{\DocumentTableSpace}{\dimexpr\ht\DocumentTableStart+\dp\DocumentTableStart
    +\LTpre+\baselineskip+#1\baselineskip*4\relax}%
  \ifdim\DocumentTableSpace>.5\textheight
    \setlength{\DocumentTableSpace}{.5\textheight}%
  \fi
  \Needspace{\DocumentTableSpace}%
}
% Measure short lists once. Small blocks move together; long ones are unboxed
% into the normal page builder and remain breakable, never scaled or clipped.
\newsavebox{\DocumentShortBlock}
\newcommand{\KeepShortBlock}[1]{%
  % A heading's page-space check must not run inside the measuring box: it
  % would inspect the outer page and embed a spurious forced break in the box.
  \par\setbox\DocumentShortBlock=\vbox{\hsize=\linewidth
    \renewcommand{\Needspace}[1]{}#1\par}%
  \ifdim\dimexpr\ht\DocumentShortBlock+\dp\DocumentShortBlock\relax<.32\textheight
    \Needspace{\dimexpr\ht\DocumentShortBlock+\dp\DocumentShortBlock+\medskipamount\relax}%
  \else
    \Needspace{5\baselineskip}%
  \fi
  % Unboxing does not add normal inter-paragraph baseline glue at the leading
  % edge. Give the measured block standard vertical separation from prior text.
  \addvspace{\medskipamount}%
  \unvbox\DocumentShortBlock
}
\preto\section{\Needspace{5\baselineskip}}
\preto\subsection{\Needspace{5\baselineskip}}
\preto\subsubsection{\Needspace{4\baselineskip}}
\clubpenalty=10000
\widowpenalty=10000
\setlist{leftmargin=1.5em,itemsep=2pt,topsep=3pt}
\providecommand{\tightlist}{\setlength{\itemsep}{1pt}\setlength{\parskip}{0pt}}
\titleformat{\section}{\large\sffamily\bfseries\color{Accent}}{}{0pt}{}
\titleformat{\subsection}{\normalsize\sffamily\bfseries\color{Accent}}{}{0pt}{}
\titleformat{\subsubsection}{\normalsize\sffamily\bfseries\color{Accent}}{}{0pt}{}
\titlespacing*{\section}{0pt}{18pt}{8pt}
\titlespacing*{\subsection}{0pt}{13pt}{6pt}
\titlespacing*{\subsubsection}{0pt}{10pt}{5pt}
\pagestyle{fancy}
\fancyhf{}
\fancyfoot[C]{\small\color{Muted}\thepage}
\renewcommand{\headrulewidth}{.35pt}
\renewcommand{\footrulewidth}{0pt}
% A single compact opening page, with profile-controlled TOC depth. Reserve a
% wider page-number column for the multi-page BRD and example editions.
\setlength{\cftbeforesecskip}{5pt}
\setlength{\cftbeforesubsecskip}{2pt}
\renewcommand{\cftsecfont}{\small\sffamily\bfseries}
\renewcommand{\cftsubsecfont}{\small\sffamily}
\renewcommand{\cftsecpagefont}{\small}
\renewcommand{\cftsubsecpagefont}{\small}
\cftsetpnumwidth{2em}
\cftsetrmarg{3em}
\newcommand{\DocumentTitle}[3]{%
  \thispagestyle{plain}
  {\sffamily\bfseries\color{Accent}\fontsize{23}{29}\selectfont #1\par}
  \vspace{3mm}
  {\sffamily\small\color{Muted}#2\par}
  \ifstrempty{#3}{}{\vspace{4mm}#3\par}%
  \begingroup
  \renewcommand{\contentsname}{Содержание}%
  \tableofcontents
  \endgroup
  \vspace{5mm}
  {\color{Accent}\rule{\linewidth}{.8pt}\par}
  \clearpage
}
% Compact standalone profile: keep the title and body on the same opening page.
% It intentionally has no table of contents or forced title-page break.
\newcommand{\CompactDocumentTitle}[3]{%
  \thispagestyle{plain}
  {\sffamily\bfseries\color{Accent}\fontsize{23}{29}\selectfont #1\par}
  \vspace{3mm}
  {\sffamily\small\color{Muted}#2\par}
  \ifstrempty{#3}{}{\vspace{4mm}#3\par}%
  \vspace{5mm}
  {\color{Accent}\rule{\linewidth}{.8pt}\par}
  \vspace{4mm}
}

\hypersetup{pdftitle={SolutionArchitectAI. Ключевые технические решения},pdfauthor={},pdfsubject={Технические решения}}
\fancyhead[L]{\small SolutionArchitectAI}
\fancyhead[R]{\small Технические решения}
\setcounter{tocdepth}{1}
\begin{document}
\DocumentTitle{SolutionArchitectAI. Ключевые технические решения}{Технические решения}{Цель документа --- зафиксировать планируемые технические решения по реализации MVP системы SolutionArchitectAI для согласования с системными архитекторами, техническими руководителями, Python- и frontend-разработчиками.

}
\section{1. Назначение и общая архитектура}\label{1-ux43dux430ux437ux43dux430ux447ux435ux43dux438ux435-ux438-ux43eux431ux449ux430ux44f-ux430ux440ux445ux438ux442ux435ux43aux442ux443ux440ux430}

SolutionArchitectAI помогает архитектору перейти от требований и корпоративного контекста к согласованной архитектуре. Рабочее место объединяет общую модель объектов и связей, её диаграммы, свойства, документы, чат с ИИ и историю решений. Предложения ИИ проходят проверку; решение об их принятии остаётся за человеком.

Frontend и backend планируется реализовать как отдельные части приложения. Frontend на \textbf{React и TypeScript} отвечает за рабочее место архитектора. Backend на \textbf{Python 3.13 и FastAPI} предоставляет HTTP API, выполняет агентный процесс, проверяет изменения и обращается к хранилищам. Браузер работает через API и не получает прямого доступа к базе данных, Git или серверу LLM.

\subsection{Локальный запуск и дальнейшее развёртывание}\label{ux43bux43eux43aux430ux43bux44cux43dux44bux439-ux437ux430ux43fux443ux441ux43a-ux438-ux434ux430ux43bux44cux43dux435ux439ux448ux435ux435-ux440ux430ux437ux432ux451ux440ux442ux44bux432ux430ux43dux438ux435}

Для локального запуска разработчиком планируется использовать \textbf{Docker Compose}. Frontend, backend и PostgreSQL запускаются в отдельных контейнерах. Отдельный контейнер \textbf{PlantUML Server} генерирует схемы внутри локального окружения. Для Git-репозиториев моделей и загруженных файлов предусматриваются сохраняемые данные. Сервер LLM подключается отдельно и не обязан входить в основной Compose.

Для будущего production планируется \textbf{Red Hat OpenShift}, где будут запускаться контейнерные образы компонентов. Compose остаётся средством локального запуска. Образы должны работать без привилегий, с конфигурацией и секретами платформы, проверками готовности и отдельным выполнением миграций. Размещение сохраняемых данных, особенно Git-репозиториев при масштабировании, требует отдельного эксплуатационного решения.

\subsection{Единая модель и представления}\label{ux435ux434ux438ux43dux430ux44f-ux43cux43eux434ux435ux43bux44c-ux438-ux43fux440ux435ux434ux441ux442ux430ux432ux43bux435ux43dux438ux44f}

Архитектурные объекты имеют устойчивые идентификаторы. Один объект может присутствовать на нескольких диаграммах: его свойства общие, а расположение элементов принадлежит конкретному представлению. Удаление элемента с диаграммы не означает удаления объекта из модели.

Интерактивное редактирование и генерация схем решают разные задачи. React Flow обслуживает работу с диаграммой в браузере. PlantUML воспроизводимо строит публикуемые схемы из проверенной модели; исходники не отправляются в публичные сервисы. Будущий обмен с Archi планируется через предметную модель и ArchiMate Model Exchange XML, независимо от внутреннего формата React Flow.

\clearpage

\KeepShortBlock{%

\section{2. Backend}\label{2-backend}

Планируется использовать следующий стек:

\begin{itemize}
\tightlist
\item
  \textbf{Python 3.13} --- язык и целевая версия backend; \textbf{uv} управляет зависимостями, окружениями и lock-файлом для воспроизводимой установки.
\item
  \textbf{FastAPI} --- HTTP API для frontend, включая получение состояния длительных операций. Типизированные контракты и OpenAPI упрощают согласование API.
\item
  \textbf{Dishka} --- контейнер внедрения зависимостей: связывает сервисы с адаптерами хранения и LLM, управляет областями их использования.
\item
  \textbf{Pydantic} --- типизированные модели входных и выходных данных и их валидация; \textbf{Pydantic Settings} загружает и проверяет конфигурацию приложения.
\item
  \textbf{SQLAlchemy} --- работа с PostgreSQL и транзакциями; \textbf{Alembic} --- управляемые миграции схемы БД. Это разделяет обычные операции с данными и изменение структуры базы.
\end{itemize}

}

Предметная модель архитектуры реализуется в Python и определяет объекты, отношения, разрешённые изменения и ограничения. \textbf{JSON используется для сериализации и восстановления модели}, сохраняемой в Git. JSON Schema проверяет формат данных, но не заменяет предметную логику.

Поддержка Python 3.13 проверена по метаданным основных пакетов и разрешению зависимостей для этой версии. Это не подтверждает поведение готового приложения; точные версии и интеграционные проверки закрепляются при реализации.

\section{3. Хранение данных}\label{3-ux445ux440ux430ux43dux435ux43dux438ux435-ux434ux430ux43dux43dux44bux445}

\textbf{PostgreSQL} хранит пользователей, проекты, задачи, чаты, сообщения, запуски агентов, состояние выполнения, согласования и журнал операций. Здесь же сохраняются контрольные точки LangGraph, необходимые для продолжения агентного процесса.

\textbf{Git} хранит историю редакций общей архитектурной модели: изменения объектов и связей, архитектурные решения и связанные представления. Рабочие ветки задач отделяют предложения от согласованной редакции; принятие изменения, архитектурное согласование и merge выполняются как отдельные действия.

\textbf{Локальное файловое или S3-совместимое хранилище} содержит загруженные документы, вложения и крупные файлы, необходимые для воспроизведения архитектурных решений. Версии источников связываются с результатами проектирования.

Эти хранилища нужны одновременно. Git даёт сравнение редакций и ветвление архитектуры, PostgreSQL --- транзакции и запросы к текущему состоянию совместной работы, файловое хранилище --- хранение крупных вложений. Один Git неудобен для сообщений и конкурентного операционного состояния. Одна PostgreSQL потребовала бы заново реализовать свойства Git; помещение крупных файлов в архитектурный репозиторий увеличивает его объём и затрудняет сопровождение.

\clearpage

\section{4. Искусственный интеллект}\label{4-ux438ux441ux43aux443ux441ux441ux442ux432ux435ux43dux43dux44bux439-ux438ux43dux442ux435ux43bux43bux435ux43aux442}

\textbf{LangChain} планируется использовать для взаимодействия с языковыми моделями. Базовые интерфейсы предоставляет \texttt{langchain-core}, а \texttt{langchain-openai} подключает настраиваемый OpenAI-совместимый сервер. Полный пакет LangChain не требуется для этого набора задач.

Единый компонент структурированной генерации на основе \textbf{LangChain и Pydantic} преобразует запросы в типизированные ответы и проверяет их структуру. Поддержка структурированного вывода зависит от модели и сервера и проверяется отдельно. Валидный ответ ещё не означает корректного архитектурного решения: изменения проходят предметные проверки и рассмотрение человеком.

\textbf{LangGraph} выполняет многоэтапный агентный процесс: извлечение требований, уточнение, проектирование и независимое ревью. Он выбран вместо собственного универсального движка выполнения, поскольку поддерживает сохранение состояния, прерывание для ответа человека и возобновление работы. \textbf{PostgreSQL checkpointer} сохраняет контрольные точки процесса; он не заменяет архитектурную модель в Git или журнал предметных операций.

\KeepShortBlock{%

\subsection{Роли агентов}\label{ux440ux43eux43bux438-ux430ux433ux435ux43dux442ux43eux432}

\begin{itemize}
\tightlist
\item
  \textbf{Агент извлечения требований} выделяет требования, противоречия и пробелы, сохраняя связь с источниками.
\item
  \textbf{Агент уточнений} формирует вопросы архитектору и учитывает полученные ответы.
\item
  \textbf{Агент проектирования} предлагает объекты, связи, решения и представления с объяснением оснований.
\item
  \textbf{Агент архитектурного ревью} независимо проверяет кандидата и сообщает о нарушениях и сомнительных предложениях.
\end{itemize}

}

Агенты работают внутри backend. Они не согласовывают собственные предложения и не выполняют merge. Изменения общих объектов показываются вместе с влиянием на связанные представления, чтобы архитектор мог принять, отклонить или вернуть предложение на уточнение.

\subsection{Локальные языковые модели}\label{ux43bux43eux43aux430ux43bux44cux43dux44bux435-ux44fux437ux44bux43aux43eux432ux44bux435-ux43cux43eux434ux435ux43bux438}

\textbf{LM Studio} планируется использовать как удобный сервер локальной модели для разработки, \textbf{vLLM} --- как вариант выделенного сервера inference. Оба подключаются отдельно через настроенный endpoint. Совпадение HTTP-формата не гарантирует одинаковых возможностей: конкретные пары модель/сервер требуют проверки структурированных ответов, потоковой выдачи и отмены. Автоматическая отправка данных внешнему провайдеру не предусматривается.

\clearpage

\section{5. Пользовательский интерфейс}\label{5-ux43fux43eux43bux44cux437ux43eux432ux430ux442ux435ux43bux44cux441ux43aux438ux439-ux438ux43dux442ux435ux440ux444ux435ux439ux441}

\textbf{React} формирует интерактивное рабочее место из компонентов; \textbf{TypeScript} проверяет типы данных и контрактов на стороне frontend. \textbf{Vite} обеспечивает локальную разработку и сборку SPA. Серверный рендеринг для MVP не требуется.

\textbf{React Router} отвечает за маршруты пространства, задачи, вида и редакции. \textbf{TanStack Query} загружает, кеширует и обновляет серверные данные FastAPI. Маршрутизация отделена от кеша; локальные несохранённые правки не подменяют подтверждённое состояние сервера.

\textbf{shadcn/ui + Tailwind CSS} выбраны для форм, таблиц, диалогов и общей темы. Компоненты доступны для изменения в проекте и сочетаются со стилизацией чата; команда сопровождает скопированный код. Обязательны светлая, тёмная и системная темы, читаемые русские подписи и управление с клавиатуры.

\textbf{assistant-ui} выбран для чата с ИИ: предоставляет компоненты сообщений и механизм работы с потоковыми ответами. Собственный адаптер соединяет его с FastAPI; история, вложения и состояние запусков сохраняются на сервере. Карточки вопросов и предложений реализуются под предметные команды приложения; облачный сервис чата не требуется.

\textbf{React Flow (\texttt{@xyflow/react})} выбран для интерактивных архитектурных диаграмм: выбор и перемещение элементов, работа со связями, масштабирование и собственные типы узлов. Приложение реализует панели свойств, правила отношений, независимое расположение элементов по видам и diff предложений ИИ с принятием или отклонением. Готовый SVG PlantUML пригоден для просмотра и публикации, но этих операций редактирования не предоставляет.

\KeepShortBlock{%

\section{6. Инструменты качества}\label{6-ux438ux43dux441ux442ux440ux443ux43cux435ux43dux442ux44b-ux43aux430ux447ux435ux441ux442ux432ux430}

\textbf{Backend:}

\begin{itemize}
\tightlist
\item
  \textbf{Ruff} --- линтинг и форматирование Python; \textbf{Pyright} --- статическая проверка типов.
\item
  \textbf{import-linter} --- контроль разрешённых зависимостей между пакетами.
\item
  \textbf{pytest} --- автоматизированные тесты предметных правил и интеграций.
\end{itemize}

}

\KeepShortBlock{%

\textbf{Frontend:}

\begin{itemize}
\tightlist
\item
  \textbf{ESLint} --- статический анализ TypeScript/React; \textbf{Prettier} --- форматирование.
\item
  \textbf{TypeScript strict} --- строгая проверка типов; \textbf{Vitest} --- модульные и компонентные тесты.
\item
  \textbf{Playwright} --- сквозные проверки работы приложения в браузере.
\end{itemize}

}

\end{document}
